\itemhead{p.47 练习 · 第4题（教材p.47）}{XJ-B1-C02-PR-P047-Y1521}
\label{XJ-B1-C02-PR-P047-Y1521}
{\footnotesize 湘教版教材p.47 / PDF第54页 / p.47 练习 / 第4题}
\field{原题}{已知二次函数 $y=ax^2+bx+c$ 有最小值 $-3$，且函数的零点为 $-1$ 和 2，求该二次函数的表达式。}
\field{参考答案}{$y=a(x+1)(x-2)$，顶点横坐标 $1/2$、最小值 $-9a/4=-3$，故 $y=\frac43x^2-\frac43x-\frac83$。}
{\small \field{关联知识}{Q01 二次函数、零点、判别式与根系关系（教材帮p.65 / PDF74）}
\field{方法与易错点}{按已知零点、顶点或系数选择交点式、顶点式或一般式；根平方和用根系关系消元。 易错：遗漏 $a
e0$；把二重根算成两个不同零点；有两正根只查积正。}
\field{答案依据}{答案册相关题组：PDF第13页（册内页12）；按当前题库原题逐项重算，不按旧版题号机械对套。}}
